\subsection{STRUCTURE INDUCED ON A SUBGROUP}

Lemma 4. Let G be a finite-dimensional Lie group, \(\Omega\) a symmetric open neighbourhood of \(e\) in G and H a subset of \(\Omega\) containing \(e\) such that the conditions \(x \in \mathrm{H}\) , \(y \in \mathrm{H}\) , \(xy^{-1} \in \Omega\) imply \(xy^{-1} \in \mathrm{H}\) . Let \(r \in \mathrm{N}_{\mathbb{K}}\) . For all \(x \in \mathrm{H}\) , let \(\mathfrak{h}_x\) be the set of \(a \in \mathrm{T}_x(\mathrm{G})\) with the following property: there exist an open neighbourhood I of 0 in K and a mapping \(f\) of class \(\mathrm{C}^r\) of I into G such that \(f(0) = x, f(I) \subset H\) , \((\mathrm{T}_0f)(1) = a\) .

\begin{enumerate}
\def\labelenumi{(\roman{enumi})}
\item
  Let \(\mathfrak{h}_e = \mathfrak{h}\) . Then \(\mathfrak{h}\) is a Lie subalgebra of \(\mathrm{L}(\mathrm{G})\) which is invariant under \(\mathrm{Ad}_{\mathrm{L}(\mathrm{G})}(\mathrm{H})\) .
\item
  \(\mathfrak{h}_x = x\mathfrak{h} = \mathfrak{h}x\) for all \(x\in \mathrm{H}\) (where \(x\mathfrak{h}\) and \(\mathfrak{h}x\) are evaluated in \(\mathrm{T(G)}\) ).
\item
  Let \(\mathbf{V}\) be a manifold of class \(\mathbf{C}^r\) , \(v_0\) a point of \(\mathbf{V}\) and \(f\) a mapping of class \(\mathbf{C}^r\) of \(\mathbf{V}\) into \(\mathbf{G}\) such that \(f(v_0) = e\) and \(f(\mathbf{V}) \subset \mathbf{H}\) . For every Lie subgroup germ \(\mathbf{H}'\) of \(\mathbf{G}\) with Lie algebra \(\mathfrak{h}, f(v) \in \mathbf{H}'\) for \(v\) sufficiently close to \(v_0\) .
\item
  For every Lie subgroup germ \(\mathbf{H}'\) of \(\mathbf{G}\) with Lie algebra \(\mathfrak{h}\) , \(\mathbf{H}' \cap \mathbf{H}\) is a neighbourhood of \(e\) in \(H'\) .
\item
  For all \(x \in \mathbf{H}\) and all \(a \in \mathfrak{b}_x\) , there exist an open neighbourhood I of 0 in K and a mapping \(f\) of class \(\mathbf{C}^{\omega}\) of I into G such that \(f(0) = x, f(I) \subset H, (T_0f)(1) = a\) . Clearly \(\mathrm{K}\mathfrak{b} = \mathfrak{b}\) and \(x\mathfrak{b}_y z = \mathfrak{b}_{xyz}\) for \(x, y, xy, xyz\) in H. This implies (ii) and the fact that \(\mathfrak{b}\) is invariant under \(\mathrm{Ad}_{\mathrm{L(G)}}(\mathrm{H})\) .
\end{enumerate}

Let \(a_1, a_2\) be in \(\mathfrak{h}\) . Let I be an open neighbourhood of 0 in K and \(f_1, f_2\) mappings of class \(\mathrm{C}^r\) of I into G such that \(f_j(0) = e, f_j(I) \subset H\) , \((\mathrm{T}_0 f_j)(1) = a_j (j = 1, 2)\) . We define \(f: I \to G\) by \(f(\lambda) = f_1(\lambda)f_2(\lambda)\) . Then \(f\) is of class \(\mathrm{C}^r\) and \(f(0) = e\) . By shrinking I if necessary, we have \(f(I) \subset H\) . On the other hand, the mapping of \(\mathrm{T}_e(G) \times \mathrm{T}_e(G)\) into \(\mathrm{T}_e(G)\) tangent to the mapping \((g, g') \to gg'\) is addition; hence \((\mathrm{T}_0 f) = a_1 + a_2\) . Hence \(a_1 + a_2 \in \mathfrak{h}\) and \(\mathfrak{h}\) is a vector subspace of L(G). Since \(x\mathfrak{h}x^{-1} = \mathfrak{h}\) for all \(x \in H\) , \((\mathrm{Ad}f_1(\lambda)).a_2 \in \mathfrak{h}\) for all \(\lambda \in I\) . The tangent mapping at 0 to the mapping \(\lambda \mapsto \mathrm{Ad}f_1(\lambda)\) is, by Proposition 44 of § 3, no. 12, the mapping \(\lambda \mapsto \mathrm{ad}(\lambda a_1)\) ; hence

\[
\left[ a _ {1}, a _ {2} \right] = (\operatorname{ad} a _ {1}). a _ {2} \in \mathfrak {h}
\]

since \(\mathfrak{h}\) is closed in \(\mathrm{L}(\mathrm{G})\) . Hence we have proved (i). In the rest of the proof we fix a Lie subgroup germ \(\mathrm{H}'\) of \(\mathrm{G}\) with Lie algebra \(\mathfrak{h}\) .

Let \(V, v_0, f\) be as in (iii). Let \(Y\) be the left foliation of \(G\) associated with \(\mathfrak{h}\) (no. 1). For all \(y \in H'\) , \(T_y(H') = y\mathfrak{h}\) . On the other hand, for all \(v \in V\) , the image of \(T_v(V)\) under \(T_v(f)\) is contained in \(\mathfrak{h}_{f(v)} = f(v)\mathfrak{h}\) (by definition of \(\mathfrak{h}_{f(v)}\) ). By Differentiable and Analytic Manifolds, R, 9.3.2, \(f\) is a morphism of \(V\) into \(Y\) . As \(H'\) is a leaf of \(Y\) (Differentiable and Analytic Manifolds, R, 9.2.8), \(f(v) \in H'\) for \(v\) sufficiently close to \(v_0\) .

Let \((a_{1},\ldots ,a_{s})\) be a basis of \(\mathfrak{h}\) . There exist an open neighbourhood I of 0 in K and mappings \(f_{1},\ldots ,f_{s}\) of class \(\mathbf{C}^r\) of I into G such that \(f_{j}(0) = e\) , \(f_{j}(\mathrm{I})\subset \mathrm{H}\) , \((\mathrm{Tf}_j)1 = a_j\) for all \(j\) . By (iii), \(f_{j}(\lambda)\in \mathrm{H}'\) for \(|\lambda |\) sufficiently small. Hence the \(f_{1}(\lambda_{1})f_{2}(\lambda_{2})\ldots f_{s}(\lambda_{s})\) constitute, for \(|\lambda_1|,\left|\lambda_2\right|,\ldots ,|\lambda_s|\) sufficiently small, a neighbourhood of \(e\) in \(\mathbf{H}'\) ; and this neighbourhood is contained in \(\mathbf{H}\) . Hence (iv).

If \(a \in \mathfrak{h}\) , there exist an open neighbourhood \(\mathbf{I}\) of 0 in \(\mathbf{K}\) and a mapping \(f\) of class \(\mathbf{C}^{\omega}\) of \(\mathbf{I}\) into \(\mathbf{G}\) such that \(f(0) = e, f(\mathbf{I}) \subset \mathbf{H}'\) , \((\mathrm{T}_0f)1 = a\) . This, with (iv), implies (v).

DEFINITION 2. \(\mathfrak{h}\) is called the subalgebra tangent to H at \(e\) .

PROPOSITION 9. Let G be a finite-dimensional Lie group and H a subgroup of G.

\begin{enumerate}
\def\labelenumi{(\roman{enumi})}
\item
  There exists on H one and only one analytic manifold structure with the following property: for all r between 1 and \(\omega\) , for every manifold V of class \(C^r\) and for every mapping f of V into H, \(f\) is of class \(C^r\) as a mapping of V into H if and only if f is of class \(C^r\) as a mapping of V into G.
\item
  With this structure, \(\mathbf{H}\) is a Lie group, the canonical injection \(i\) of \(\mathbf{H}\) into \(\mathbf{G}\) is an immersion and \(\mathbf{L}(i)(\mathbf{L}(\mathbf{H}))\) is the Lie subalgebra tangent at \(e\) to \(\mathbf{H}\) .
\end{enumerate}

In (i) the uniqueness is obvious. We prove the existence. Let \(\mathfrak{b}\) be the Lie algebra tangent at \(e\) to H. Let H' be a Lie subgroup germ of G with Lie algebra \(\mathfrak{b}\) . By replacing H' by an open subgroup germ of H', it can be assumed that H' ⊆ H (Lemma 4 (iv)). For all \(x \in \mathrm{H}\) , \(xH'x^{-1}\) is a Lie subgroup germ of G with Lie algebra \(x \mathfrak{h}x^{-1} = \mathfrak{b}\) . Hence H' ∩ \((xH'x^{-1})\) is open in H' (no. 2, Theorem 3) and the mapping y → \(xyx^{-1}\) is an isomorphism of H' ∩ \(x^{-1}H'x\) onto \(xH'x^{-1}\) ∩ H'. Using Proposition 18 of § 1, no. 9, there exist an open Lie subgroup germ W of H' and a Lie group structure on H with the following properties: W is open in H and the manifold structures on H and H' induce the same structure on W. It then follows that the canonical injection i of H into G is an immersion and that L(i)(L(H)) = L(H') = \(\mathfrak{b}\) . Moreover, let V and f be as in (i). If f: V → H is of class \(C^r\), \(i \circ f\): V → G is of class \(C^r\). Suppose that \(i \circ f\): V → G is of class \(C^r\); then we prove that f: V → H is of class \(C^r\). By translation, it suffices to consider the case where there exists v0 ∈ V such that f(v0) = e and to prove that f: V → H is of class \(C^r\) in an open neighbourhood of v0. Now, by Lemma 4 (iii), f(v) ∈ H' for v sufficiently close to v0, whence our assertion. Thus we have proved (i) and (ii) has been obtained on the way.

DEFINITION 3. The Lie group structure on H defined in Proposition 9 is called the structure induced on H by the Lie group structure on G.

If H is a Lie subgroup of G, its Lie group structure is induced by that on G (Differentiable and Analytic Manifolds, R, 5.8.5).

If \(\mathbf{G} = \mathbf{R}\) and \(\mathbf{H} = \mathbf{Q}\) , then \(\mathfrak{h} = \{0\}\) and hence the structure induced on \(\mathbf{H}\) is the discrete Lie group structure. Similarly if \(\mathbf{G} = \mathbf{C}\) (considered as a complex Lie group) and \(\mathbf{H} = \mathbf{R}\) .
% semantic-fragments: legacy-input-seam
\relax{}
